\documentclass[11pt,a4paper]{article}

\usepackage[margin=2.55cm]{geometry}
\usepackage{amsmath,amssymb,amsthm,mathtools}
\usepackage{booktabs}
\usepackage{microtype}
\usepackage[hidelinks]{hyperref}
\usepackage{enumitem}

\newtheorem{definition}{Definition}[section]
\newtheorem{proposition}[definition]{Proposition}
\newtheorem{remark}[definition]{Remark}

\newcommand{\exec}{\mathrm{exec}}
\newcommand{\pre}{\mathrm{pre}}
\newcommand{\add}{\mathrm{add}}
\newcommand{\del}{\mathrm{del}}
\newcommand{\clamp}{\operatorname{clamp}}
\newcommand{\softmax}{\operatorname{softmax}}
\newcommand{\LSE}{\operatorname{LSE}}
\newcommand{\pos}{\operatorname{pos}}
\newcommand{\NN}{\mathbb{N}}
\newcommand{\RR}{\mathbb{R}}
\newcommand{\cK}{\mathcal{K}}
\newcommand{\cT}{\mathcal{T}}
\newcommand{\cC}{\mathcal{C}}
\newcommand{\cR}{\mathcal{R}}
\newcommand{\cS}{\mathcal{S}}
\newcommand{\cP}{\mathcal{P}}

\title{The PSPR Planning Algorithm\\[2pt]
\large Constraint Partitioning, Relaxed Penalty Merging, and Hierarchical Refinement}
\author{PSPR project}
\date{\today}

\begin{document}
\maketitle

\section{The planning problem}

A grounded hierarchical planning problem is a tuple
\begin{equation}
P=\langle F,\,A,\,C,\,M,\,s_0,\,TN_0,\,g\rangle,
\end{equation}
where $F$ is a finite set of ground atoms, $A$ a set of primitive action schemas,
$C$ a set of compound task schemas, $M$ a set of methods,
$s_0\subseteq F$ the initial state, $TN_0=\langle T_0,\prec_0,\approx_0\rangle$
the initial task network, and $g$ a state-goal formula, possibly absent.
For a primitive action $a$ with precondition $\pre_a\subseteq F$,
add effects $\add_a\subseteq F$, and delete effects $\del_a\subseteq F$,
the state transition is
\begin{equation}
\tau_a(s)=
\begin{cases}
(s\setminus\del_a)\cup\add_a & \text{if } s\models\pre_a,\\
\text{undefined} & \text{otherwise}.
\end{cases}
\end{equation}
A task network is $TN=\langle T,\prec,\approx\rangle$: a finite multiset of task
occurrences, an acyclic ordering $\prec\;\subseteq T\times T$, and equality and
inequality constraints $\approx$. A method
\begin{equation}
m=\langle c(\vec{x}),\,\phi_m,\,TN_m\rangle
\end{equation}
decomposes the compound task $c(\vec{x})$ into the subtask network $TN_m$.
A decomposition step replaces an occurrence $t=c(\vec{t})$ by a grounded method
instance $m\sigma$ and updates the ordering by
\begin{equation}
\prec'=\prec\;\cup\;\prec_m\sigma
\;\cup\;\{(u,v)\mid u\prec t,\;v\in T_m\sigma\}
\;\cup\;\{(v,w)\mid v\in T_m\sigma,\;t\prec w\}.
\end{equation}
A decomposition tree $D$ is primitive when every task in its frontier $L(D)$
is a primitive action. A plan $\pi=\langle a_1,\dots,a_k\rangle$ is a
linearisation of $L(D)$ that respects $\prec_D$ and $\approx_D$; it is
executable from $s_0$ when each $a_i$ is applicable in the state produced by
its predecessors, and it is a solution when, in addition, the final state
satisfies $g$.

\paragraph{Notation.}
For a task occurrence $t$ we write $\mathrm{head}(t)$ for its task symbol and
$\mathrm{args}(t)$ for its ground arguments. The projection of a set of ground
atoms onto a subset $X\subseteq F$ is $s[X]=\{a\in s\mid a\in X\}$.

\section{The abstraction computed by the evaluation tree}

The evaluation tree does not solve $P$ directly. It solves a
\emph{top-level projection} $\alpha(P)$ in which compound tasks are opaque.
For every grounded task occurrence $t\in T_0$ introduce a marker atom
$\exec_t\notin F$. A primitive occurrence $t$ with $\mathrm{head}(t)=a$ is
compiled into the operator
\begin{equation}
\pre(o_t)=\pre_a\;\wedge\;\bigwedge_{u\prec_0 t}\exec_u,
\qquad
\add(o_t)=\add_a\cup\{\exec_t\},
\qquad
\del(o_t)=\del_a,
\label{eq:abstract-action}
\end{equation}
and a compound occurrence $t$ with method instances $M(t)$ gets one operator
per instance $m\sigma$:
\begin{equation}
\pre(o_{t,m})=\phi_m\sigma\;\wedge\;\bigwedge_{u\prec_0 t}\exec_u,
\qquad
\add(o_{t,m})=\{\exec_t\},
\qquad
\del(o_{t,m})=\varnothing.
\label{eq:abstract-method}
\end{equation}
The ordering of $TN_0$ is thus enforced operationally by predecessor markers,
not by a separate constraint layer. The state goal $g$ is not part of the
abstraction; it is checked after planning, either by the strict validator or by
the refinement stage of Section~\ref{sec:refinement}.

\begin{remark}[The liftability gap]
The projection is neither a sound relaxation nor a complete abstraction of $P$.
It can \emph{reject} a valid HTN plan: if the subtasks of a method for $t_1$
establish an atom $p$ that is a precondition of a method for a later task $t_2$,
then in $\alpha(P)$ the operator for $t_2$'s method cannot fire, because
$t_1$ only writes $\exec_{t_1}$. Conversely, it can \emph{accept} a top-level
plan that no decomposition supports: an abstract method operator only tests
$\phi_m$, while the first subtask of $m\sigma$ may be inapplicable in $s_0$.
No inclusion holds in general between the top-level plans of $\alpha(P)$ and
the top-level plans induced by HTN solutions of $P$. This gap is why PSPR
offers the bounded refinement of Section~\ref{sec:refinement}; the default
path returns a plan of the projection plus an independent validation report.
\end{remark}

\paragraph{Detecting the gap statically.}
For a fixed grounding the first half of the gap is decidable by a may-write
analysis. Let $W(t)$ be the set of ground literals that some decomposition of
the top-level task $t$ can write, computed as the least fixpoint of
\begin{equation}
W(t)=\bigcup_{m\in M(t)}\;\bigcup_{s\in TN_m}W(s)
\qquad\text{and}\qquad
$W(a)=\bigcup_{\alpha\in\mathrm{eff}(a)}\{(p,v):(p\mathrel{:=}v)\in\alpha\}$,
\end{equation}
where a primitive task $a$ uses its effect alternatives and a compound task
$t$ uses its grounded methods. The least fixpoint converges on recursive
methods and contains exactly the atoms that a finite decomposition can write.
Let $R(t)$ be the literals of the DNF of the preconditions of every grounded
method of $t$ (of the action schemas when $t$ is primitive), together with,
for each literal, the number of clauses that contain it. For every ordered
pair $t_1\prec^{+}t_2$ and every $p\in W(t_1)\cap R(t_2)$ the pair is a
candidate gap. A sign match is a \emph{hidden satisfier} (the projection cannot
see the establishment and may disable $t_2$'s operator); a sign mismatch is a
\emph{hidden falsifier} (the projection keeps believing the literal). The
implementation is \texttt{evaluation\_tree::liftability}. It over-approximates
on purpose -- preconditions are ignored, disjunctive effects are merged, and
nested subtasks are expanded over every method grounding -- so a reported pair
is a candidate rather than a proof, and a clean report is not a proof of
liftability. Each report entry records the producer method, the consumer method
or action, the witness subtask, whether the literal occurs in every DNF clause
of the consumer, whether it already holds initially, and whether a primitive
top-level task can write it anyway.

\paragraph{Repairing the gap with implicit orderings.}
The same analysis drives a repair. Among the task pairs that the network
leaves unordered in both directions, take every pair where a method of $t_1$
establishes a literal that every clause of a precondition of $t_2$ requires,
the literal does not hold initially, and no primitive top-level task can write
it; add the edge $t_1\prec t_2$. A hidden falsifier is handled symmetrically:
if $t_1$ destroys a literal that $t_2$ requires and that holds initially, the
edge $t_2\prec t_1$ is added instead. Candidates that would close a cycle are
dropped, and the remaining edges are inserted into the task network before
clustering. They are then ordinary ordering constraints: the incidence graph
connects their endpoints, the leaves close the ordering transitively, and the
merge nodes validate them exactly. The abstract plan therefore executes the
producer before the consumer, and the refinement stage threads the established
atom into the consumer's method precondition. The hierarchical solver retries
with this repair when the first attempt fails and the analysis finds an edge
(\texttt{EvaluationOptions::implicit\_orderings}, \texttt{--repair-orderings}
on the command line).

\section{Grounding and the task/constraint incidence graph}

\subsection{Grounding}
The free parameters of $TN_0$ are enumerated over the typed constants and
objects of the problem. A parameter ranges over the symbols whose declared type
is a subtype of the parameter type. If the enumeration produces several
groundings, their task lists are concatenated and each grounding contributes
its own copy of the ordering and equality constraints. The number of
assignments is bounded by the grounding budget $B_g$
(\texttt{EvaluationOptions::max\_groundings}, default $256$ in the CLI).

Action and method formulas are compiled in four steps:
\begin{enumerate}[leftmargin=1.6em,itemsep=1pt]
\item quantifiers $\exists$ and $\forall$ are expanded over the typed universe;
\item $\mathit{Imply}$, $\mathit{Xor}$ and composite $\mathit{Not}$ are
      rewritten into negation normal form over literals;
\item $\mathit{Equals}$ is evaluated against the current bindings;
\item preconditions are converted to disjunctive normal form (DNF), and effects
      are compiled into conditional effects. Every DNF clause and every
      disjunctive effect alternative becomes one operator variant.
\end{enumerate}
The DNF expansion is capped at $4096$ clauses or alternatives per formula.
$\mathit{Probabilistic}$ is the only constructor that cannot be represented;
it is counted in \texttt{EvaluationTree::unmodeled} and makes leaf construction
fail.

\subsection{The incidence graph}
Let $\cT=\{t_1,\dots,t_{n_1}\}$ be the grounded task occurrences and let
$\cK=\{c_1,\dots,c_{n_2}\}$ be the constraints they touch. A constraint node is
one of:
\begin{itemize}[leftmargin=1.6em,itemsep=1pt]
\item an atom occurring in a matching action schema (precondition or effect) or
      in a matching method precondition (and, when summaries are enabled, in a
      summary effect);
\item an ordering pair $(t_i,t_j)$ of the network;
\item an equality or inequality constraint between two ground terms.
\end{itemize}
The incidence graph is the bipartite graph with adjacency matrix
\begin{equation}
A_{ij}=
\begin{cases}
1 & \text{if task }t_i\text{ touches constraint }c_j,\\
0 & \text{otherwise}.
\end{cases}
\end{equation}
Its connected components are used as \emph{required groupings}: independent
parts of the network are clustered separately and never mixed by the spectral
step. Explicit groupings can be supplied through
\texttt{EvaluationOptions::groupings}; tasks not mentioned in any grouping
become singleton groups.

\section{Spectral clustering of tasks}
\label{sec:clustering}

Each grouping is clustered independently on its induced subgraph. Let the
group have $n_1$ type-I nodes, $n_2$ type-II nodes, $m$ edges and density
\begin{equation}
p=\frac{m}{n_1n_2}.
\end{equation}
Two regimes are distinguished by the signal condition
$n_1n_2p^2\ge\log n_1$.

\subsection{Dense regime: \texorpdfstring{$\mathrm{Spec}$}{Spec} and
\texorpdfstring{$\mathrm{AdaSpec}$}{AdaSpec}}
Let
\begin{equation}
S=H(AA^\top)
\end{equation}
be the \emph{hollowed} Gram matrix: $S_{ij}=(AA^\top)_{ij}$ for $i\ne j$ and
$S_{ii}=0$. Hollowing is essential because the diagonal of $AA^\top$ is of
order $n_2p_{\max}$, while the population diagonal of $PP^\top$ is of order
$n_2p_{\max}^2$; keeping it would dominate the spectrum. Let
$S=U\Lambda U^\top$ with
$\lambda_1\ge\lambda_2\ge\dots\ge\lambda_{n_1}$ and orthonormal columns
$u_1,\dots,u_{n_1}$.

$\mathrm{AdaSpec}$ selects the embedding rank
\begin{equation}
r=\max\left\{i\in\{1,\dots,n_1-1\}:\;
\lambda_i-\lambda_{i+1}>\theta\right\},
\qquad
\theta=\frac{n_1n_2p^2}{\log\log n_1},
\label{eq:adaspec-rank}
\end{equation}
where the threshold is replaced by $n_1n_2p^2$ when $\log\log n_1\le0$ or is
non-finite. If no gap exceeds $\theta$, the implementation falls back to the
largest gap, provided it exceeds $10^{-9}\max(|\lambda_1|,1)$; otherwise the
group stays a single cluster. The embedding of task $i$ is
\begin{equation}
u(i)=\bigl(U_{i,1},\dots,U_{i,r}\bigr)\in\RR^r,
\end{equation}
and the number of communities is $k=r$.

\subsection{Sparse regime: normalized Laplacian}
When $n_1n_2p^2<\log n_1$ the hollowed Gram matrix is too noisy. The
implementation uses the symmetric normalized Laplacian of the similarity graph
$S$:
\begin{equation}
L=I-D^{-1/2}SD^{-1/2},
\qquad
D_{ii}=\sum_{j}S_{ij},
\end{equation}
with spectrum $0=\mu_1\le\mu_2\le\dots\le\mu_{n_1}$. The number of communities
is the classic eigengap estimate
\begin{equation}
k=\arg\max_{c\ge1}\bigl(\mu_{c+1}-\mu_c\bigr),
\end{equation}
and the embedding is given by the eigenvectors of
$\mu_2,\dots,\mu_k$.

\subsection{Assignment and clusters}
Both regimes partition the embedded rows with $k$-medians:
\begin{equation}
\min_{\mu_1,\dots,\mu_k}\;\sum_{i=1}^{n_1}\bigl\|u(i)-\mu_{\ell(i)}\bigr\|_2,
\end{equation}
where rows are assigned to the nearest centre in the Euclidean norm and each
centre is updated to the coordinate-wise median of its cell. The iteration is
deterministic, seeded with SplitMix64. Cluster $j$ consists of its type-I
tasks $\cT_j$ together with every type-II node adjacent to at least one of
them:
\begin{equation}
\cK_j=\bigl\{c\in\cK:\exists t\in\cT_j,\;A_{t,c}=1\bigr\}.
\end{equation}
A constraint is \emph{shared} when it belongs to more than one $\cK_j$:
\begin{equation}
\cS=\Bigl\{c\in\cK:\bigl|\{j:c\in\cK_j\}\bigr|>1\Bigr\},
\end{equation}
and \emph{local} to cluster $j$ otherwise. Local constraints are solved inside
the leaf; shared constraints are delegated to the merge nodes.

\section{Agglomeration and the evaluation tree}
\label{sec:agglomeration}

\subsection{Pairwise distance}
For two clusters $j$ and $j'$ the implementation uses the Jaccard distance
\begin{equation}
d(j,j')=1-\frac{|\cK_j\cap\cK_{j'}|}{|\cK_j\cup\cK_{j'}|},
\end{equation}
with $d(j,j')=0$ when both constraint sets are empty.

\subsection{The merge-level ILP}
Let $k$ be the number of clusters and $L$ the number of levels
($L=k-1$ by default). The hierarchy is represented by integer merge-level
variables
\begin{equation}
M_{ab}\in\{1,\dots,L\},\qquad a\ne b,
\end{equation}
where $M_{ab}$ is the level at which $a$ and $b$ first belong to the same
cluster. For every ordered triple $(a,b,c)$ of distinct clusters the
formulation introduces binary indicators $O_{abc}$, $Z^{ab}_{abc}$ and
$Z^{bc}_{abc}$, and the constraints
\begin{align}
-L\;\le\;&M_{ac}-M_{ab}-(L+1)O_{abc}\;\le\;0,
\label{eq:gilpin-o}\\
-L\;\le\;&M_{ab}-M_{ac}-(L+1)Z^{ab}_{abc}+1\;\le\;0,
\label{eq:gilpin-zab}\\
-L\;\le\;&M_{bc}-M_{ac}-(L+1)Z^{bc}_{abc}+1\;\le\;0,
\label{eq:gilpin-zbc}\\
&Z^{ab}_{abc}+Z^{bc}_{abc}\;\ge\;1.
\label{eq:gilpin-z}
\end{align}
Constraint \eqref{eq:gilpin-o} forces $O_{abc}=1$ exactly when
$M_{ac}<M_{ab}$, i.e.\ when $b$ joins $a$ before $c$ does.
Constraints \eqref{eq:gilpin-zab}--\eqref{eq:gilpin-z} encode the dendrogram
validity condition that $M_{ac}$ cannot be strictly larger than both $M_{ab}$
and $M_{bc}$; equivalently, none of the three points may join the other two
last. The clustering objective is
\begin{equation}
Q(A)=\sum_{a}\sum_{b\ne a}\sum_{c\ne a,b}
\bigl(d(a,c)-d(a,b)\bigr)\,O_{abc},
\label{eq:Q}
\end{equation}
which is the sum, over triples where $a$ and $b$ merge before $a$ and $c$, of
the gain obtained by placing the closer neighbour $b$ first.

\subsection{Constraint locality}
For every shared constraint $c$ with cluster set
$S_c=\{j:c\in\cK_j\}$ the formulation adds a continuous height variable
$h_c\in[1,L]$ and the constraints
\begin{equation}
h_c\;\ge\;M_{ab}\qquad\text{for all distinct }a,b\in S_c.
\end{equation}
Since the objective penalizes $\sum_c h_c$, at optimality
\begin{equation}
h_c=\max_{a,b\in S_c}M_{ab},
\end{equation}
which is exactly the level of the lowest common ancestor of the clusters that
share $c$. The complete hierarchy objective is therefore
\begin{equation}
\max_{A}\;\;Q(A)-\gamma\,G(A),
\qquad
G(A)=\sum_{c\in\cS}h_c,
\label{eq:hierarchy-objective}
\end{equation}
where $\gamma$ is \texttt{Options::global\_weight} (default $1$). A constraint
handled by the first merge costs $1$; one pushed to the root costs $L$. The
penalty keeps shared constraints low in the tree, where they are checked
against smaller subplans.

\subsection{From merge levels to a binary dendrogram}
At every level $\ell$ the clusters are the connected components of the relation
``$M_{ab}\le\ell$'', computed with a union-find structure. A rooted binary
dendrogram is built by walking the levels from $1$ to $L$: each non-singleton
cluster at level $\ell$ is formed from the clusters of level $\ell-1$ that it
contains, and a multi-way merge is binarised by folding its children in
ascending order. Leaves are the $k$ clusters; the $k-1$ internal nodes are
created in chronological order, and internal node $k+\ell$ receives evaluation
id $\ell$.

Each shared constraint is attached to the \emph{first} internal node, in
creation order, whose member set contains every cluster in $S_c$. Since nodes
are created bottom-up, this is the lowest common ancestor of $S_c$. The merge
node therefore receives exactly the constraints whose clusters first meet
there.

\subsection{The scheduling ILP}
Let $I$ be the set of internal nodes, $N=|I|=k-1$, and let $E_A$ be the
parent--child edges of the dendrogram, directed from child to parent. Binary
variables $x_{v,t}$ indicate that internal node $v$ is evaluated at time $t$:
\begin{align}
\sum_{t=1}^{N}x_{v,t}&=1 && \forall v\in I,
\label{eq:sched-once}\\
\sum_{v\in I}x_{v,t}&=1 && \forall t\in\{1,\dots,N\},
\label{eq:sched-slot}
\end{align}
with evaluation position $p_v=\sum_{t=1}^{N}t\,x_{v,t}$. Every child is
evaluated before its parent, and an optional evaluation DAG
$B=(V_B,E_B)$ adds precedences:
\begin{align}
p_u+1&\le p_v && \forall (u,v)\in E_A,\;u,v\in I,
\label{eq:sched-tree}\\
p_{f(a)}+1&\le p_{f(b)} && \forall (a,b)\in E_B,
\label{eq:sched-dag}
\end{align}
where $f$ maps DAG vertices to internal nodes. The ILP minimizes
\begin{equation}
C(A,\sigma)=\sum_{v\in I}p_v.
\end{equation}
The union of \eqref{eq:sched-tree} and \eqref{eq:sched-dag} must be acyclic;
otherwise the schedule is rejected. The resulting order
$\sigma$ is a topological order of the merge nodes.

\subsection{Joint objective}
The selected tree and schedule are scored by
\begin{equation}
Q(A)-\lambda\,C(A,\sigma)-\gamma\,G(A),
\label{eq:joint}
\end{equation}
where $\lambda$ is \texttt{Options::lambda} (default $0$). Since every complete
binary hierarchy over $k$ leaves has exactly $k-1$ internal nodes, $C$ is
nearly constant across hierarchies and the ILP of
\eqref{eq:hierarchy-objective} dominates the choice. An optional
planning-style search over partial dendrograms can be enabled
(\texttt{use\_search}); it reuses the generic best-first engine and picks its
candidate only when it improves \eqref{eq:joint}. In the default configuration
only the exact ILP is solved.

\section{Leaf subproblems}
\label{sec:leaves}

For cluster $j$ let $L_j$ be the set of atoms in its local constraints. The
leaf planning task has binary variables
\begin{equation}
V_j=\{x_a:a\in L_j\}\;\cup\;\{e_t:t\in\cT_j\},
\end{equation}
where $x_a$ is the truth value of atom $a$ and $e_t$ is the execution marker
of task $t$. The initial assignment is
\begin{equation}
x_a=1\iff a\in s_0,
\qquad
e_t=0,
\end{equation}
and the goal is
\begin{equation}
G_j=\{e_t=1:t\in\cT_j\}.
\end{equation}
Shared atoms are dropped from the leaf entirely: they neither appear as
preconditions nor as effects of leaf operators. This is what makes the leaves
independent; the missing shared atoms are the interface that the merge must
repair.

\subsection{Ordering closure}
Let $\prec_j$ be the ordering edges induced by the local ordering constraints
between tasks of $\cT_j$. The implementation computes the transitive closure
$\prec_j^*$ (Floyd--Warshall over the cluster positions) and rejects a cyclic
closure. Every operator for task $t$ contains the predecessor markers
\begin{equation}
\prec_j^*(t)=\{e_u=1:u\prec_j^*t\}.
\end{equation}

\subsection{Action variants}
Let $\phi$ be a precondition in DNF,
$\phi=\bigvee_{\kappa\in K}\bigwedge_{\ell\in\kappa}\ell$, and let
$\alpha=\{(\gamma,a,v)\}$ be one deterministic effect alternative, where
$\gamma$ is a condition, $a$ an atom and $v\in\{0,1\}$ its new value. For every
pair $(\kappa,\alpha)$ the leaf contains one operator
\begin{equation}
\begin{aligned}
\pre(o)&=\prec_j^*(t)\;\cup\;
\bigl\{x_a=\kappa(a):a\in\kappa\cap L_j\bigr\},\\
\mathrm{eff}(o)&=\bigl\{x_a:=v\text{ if }\gamma:(a,v)\in\alpha,\;a\in L_j\bigr\}
\cup\{e_t:=1\}.
\end{aligned}
\label{eq:leaf-action}
\end{equation}
A task that names an action is compiled through its action schemas; methods are
used only when no action schema matches. The restriction to $L_j$ in the
condition and effect sets is intentional: constraints outside the leaf belong
to a merge node. Effect atoms are still recorded in the leaf's
\texttt{effect\_atoms}, so that the merge can check which shared atoms a child
plan establishes.

\subsection{Compound tasks}
A compound task $t$ with no matching action schema gets one operator per
grounded method instance $m\sigma$ and per DNF clause $\kappa$ of
$\phi_m\sigma$:
\begin{equation}
\pre(o)=\prec_j^*(t)\;\cup\;
\bigl\{x_a=\kappa(a):a\in\kappa\cap L_j\bigr\},
\qquad
\mathrm{eff}(o)=\{e_t:=1\}.
\label{eq:leaf-method}
\end{equation}
If the domain declares no method whose task name and arity match $t$, a
synthetic operator \texttt{task[$t$]} with no effect other than $e_t:=1$ is
used. When primitive summaries are enabled
(Section~\ref{sec:refinement}), a summarisable method is replaced by its macro
operator; methods outside the summarisable fragment stay opaque.

\subsection{Heuristic search}
The leaf task is a finite multi-valued planning task. It is solved by greedy
best-first search with the FF heuristic. The relaxed planning graph assigns
each fact $p$ a level
\begin{equation}
h(p)=\min_{o:\,p\in\add_o}\Bigl(\max_{q\in\pre_o}h(q)+\mathrm{cost}(o)\Bigr),
\end{equation}
starting from $h(p)=0$ for facts of the current state; axioms are applied to a
fixed point on every layer. If a goal fact has infinite level the state is
dead. Otherwise $h_{\mathrm{FF}}$ is the number of operators in the relaxed
plan obtained by regressing from the goal facts through their cheapest
achievers; the relaxed-plan operators that are applicable in the real state
are returned as \emph{preferred} actions. Search maintains two open lists per
heuristic (preferred and non-preferred), uses a closed list, evaluates
heuristics lazily when a node is expanded, and expands at most
$10^5$ nodes per leaf. Every operator has unit cost, so the leaf returns the
first plan found, not a cost-optimal one.

\section{Merging as a relaxed MINLP}
\label{sec:merge}

\subsection{The SGPlan penalty formulation}
Let a merge node $v$ have $n_c$ children $C_1,\dots,C_{n_c}$ with plans
\begin{equation}
\pi_i=\langle o_{i,1},\dots,o_{i,l_i}\rangle,
\qquad
\pi_i\models P_i,
\end{equation}
and let $\cS_v$ be the shared constraints attached to $v$. SGPlan's
constraint-partitioning formulation is the mixed-integer nonlinear program
\begin{equation}
\begin{aligned}
\min_{\pi_1,\dots,\pi_{n_c},\;w}\quad &
\sum_{i=1}^{n_c}\mathrm{cost}_i(\pi_i)
+\sum_{c\in\cS_v}w_c\,v_c(\pi_1,\dots,\pi_{n_c})\\
\text{s.t.}\quad &
v_c(\pi_1,\dots,\pi_{n_c})=0 && \forall c\in\cS_v,\\
& \pi_i\in\mathrm{Plans}(P_i) && \forall i,\\
& w_c\ge0 && \forall c\in\cS_v,
\end{aligned}
\label{eq:minlp}
\end{equation}
where $v_c$ is the violation of shared constraint $c$ and $w_c$ its penalty
weight. PSPR does not re-plan inside the children: it fixes the $\pi_i$ and
optimizes only how they are interleaved.

\subsection{Continuous relaxation}
Let
\begin{equation}
L=\sum_{i=1}^{n_c}l_i,
\qquad
\ell_i=\frac{l_i}{L},
\qquad
q_i\in[0,1]
\end{equation}
be the total number of steps, the normalized length of child $i$, and its
continuous start position. Step $k$ of child $i$ (with $k=0,\dots,l_i-1$)
occurs at the normalized time
\begin{equation}
\tau_{i,k}=q_i+\frac{k}{L},
\end{equation}
and the child completes at $q_i+\ell_i$. The makespan is replaced by a
log-sum-exp smooth maximum
\begin{equation}
M_T(q)=T\log\sum_{i=1}^{n_c}
\exp\!\left(\frac{q_i+\ell_i}{T}\right),
\label{eq:smoothmax}
\end{equation}
whose gradient with respect to $q_i$ is the softmax weight
\begin{equation}
\frac{\partial M_T}{\partial q_i}
=
\frac{\exp\!\bigl((q_i+\ell_i)/T\bigr)}
{\sum_{j}\exp\!\bigl((q_j+\ell_j)/T\bigr)}.
\end{equation}

For a shared ordering constraint $a\prec b$, where $a$ is step $k$ of child
$i$ and $b$ is step $l$ of child $j$, define the signed margin
\begin{equation}
g_{ab}(q)=q_j+\frac{l}{L}-q_i-\frac{k}{L}-\frac{1}{L}.
\label{eq:margin}
\end{equation}
It is non-negative exactly when $a$ is scheduled strictly before $b$ with at
least one step of separation. Its exact violation is $[-g_{ab}]_+$, and its
smooth violation is the one-sided hinge
\begin{equation}
v_{ab}(q)=T\,\mathrm{softplus}\!\left(-\frac{g_{ab}(q)}{T}\right)
=T\log\!\left(1+\exp\!\left(-\frac{g_{ab}(q)}{T}\right)\right),
\label{eq:hinge}
\end{equation}
with derivative
\begin{equation}
\frac{\partial v_{ab}}{\partial g_{ab}}
=-\sigma\!\left(-\frac{g_{ab}}{T}\right),
\qquad
\sigma(z)=\frac{1}{1+e^{-z}}.
\end{equation}
Hence the penalty stops pushing as soon as the constraint is comfortably
satisfied.

Atom and equality constraints do not depend on the schedule. They enter the
objective as constants, where an unsatisfied atom or equality contributes
$w_c\,\phi(1)$ and a satisfied one contributes $0$; they still drive
feasibility and the weight adaptation.

\subsection{The relaxed objective}
With penalty function $\phi$, make-span weight $c$, and weights $w_c$, the
optimized objective is
\begin{equation}
J_T(q;w)=c\,M_T(q)+\sum_{c\in\cS_v}w_c\,\phi\bigl(v_c(q)\bigr),
\label{eq:J}
\end{equation}
and its gradient is
\begin{equation}
\begin{aligned}
\frac{\partial J_T}{\partial q_i}
&=
c\,\frac{\exp\bigl((q_i+\ell_i)/T\bigr)}
{\sum_{j=1}^{n_c}\exp\bigl((q_j+\ell_j)/T\bigr)}\\
&\quad+
\sum_{\substack{c=(a,b)\in\cS_v\\ b\text{ in child }i}}
w_c\,\phi'(v_c)\left(-\sigma\!\left(-\frac{g_c}{T}\right)\right)\\
&\quad-
\sum_{\substack{c=(a,b)\in\cS_v\\ a\text{ in child }i}}
w_c\,\phi'(v_c)\left(-\sigma\!\left(-\frac{g_c}{T}\right)\right).
\end{aligned}
\label{eq:gradient}
\end{equation}
The three penalty spaces are
\begin{equation}
\phi_{\mathrm{lin}}(v)=v,
\qquad
\phi_{\mathrm{quad}}(v)=v^2+v,
\qquad
\phi_{\mathrm{exp}}(v)=\frac{e^{4v}-1}{4},
\end{equation}
with slopes $1$, $2v+1$ and $e^{4v}$.

\subsection{Adam with box projection}
The relaxed problem is minimized by Adam. With gradient $g_t$, moments
\begin{equation}
m_t=\beta_1m_{t-1}+(1-\beta_1)g_t,
\qquad
v_t=\beta_2v_{t-1}+(1-\beta_2)g_t^2,
\end{equation}
bias corrections $\hat m_t=m_t/(1-\beta_1^t)$ and
$\hat v_t=v_t/(1-\beta_2^t)$, the update is
\begin{equation}
q^{(t+1)}=\Pi_{[0,1]}\!\left(q^{(t)}
-\eta\,\frac{\hat m_t}{\sqrt{\hat v_t}+\epsilon}\right),
\label{eq:adam}
\end{equation}
where $\Pi_{[0,1]}$ is the coordinate-wise projection onto the unit box. The
defaults are $\eta=0.05$, $\beta_1=0.9$, $\beta_2=0.999$, $\epsilon=10^{-8}$,
temperature $T=0.05$, make-span weight $c=1$, $10$ outer rounds of $200$ Adam
steps each.

\subsection{Adaptive weights and penalty continuation}
After every outer round the exact violations $v_c$ are measured. The weights
are updated multiplicatively,
\begin{equation}
w_c\leftarrow
\begin{cases}
\clamp\bigl(\gamma\,w_c,\;w_{\min},\;w_{\max}\bigr)
& \text{if } v_c>\tau,\\[2pt]
\clamp\bigl(\delta\,w_c,\;w_{\min},\;w_{\max}\bigr)
& \text{if } v_c\le\tau,
\end{cases}
\label{eq:weights}
\end{equation}
with growth $\gamma=4$, decay $\delta=1/2$, floor $w_{\min}=10^{-3}$, ceiling
$w_{\max}=10^{9}$ and tolerance $\tau=10^{-3}$. Decaying satisfied
constraints keeps the active set well conditioned.

Let $V^{(r)}$ be the total exact violation after round $r$. A round is
\emph{stalled} when
\begin{equation}
V^{(r)}>(1-\rho)\,V^{(r-1)},
\qquad \rho=0.05.
\end{equation}
A stall advances the continuation from $\phi_{\mathrm{lin}}$ to
$\phi_{\mathrm{quad}}$ and then to $\phi_{\mathrm{exp}}$, and resets Adam's
moments so that the new curvature is measured from scratch. The best iterate
of the whole run, ordered lexicographically by (total violation, objective),
is retained, so a late divergent round cannot degrade the result.

If the merge node has no ordering constraints, if \texttt{iterations} is zero,
or if the initial total violation is already at most $\tau$, Adam is skipped
and the initial (or learned) positions are decoded directly.

\subsection{Decoding and exact validation}
The relaxed positions are decoded by greedy list scheduling. Let
$k_i\leftarrow0$ for every child. Repeatedly select
\begin{equation}
i^*=\arg\min_{\,i:\,k_i<l_i}\left(q_i+\frac{k_i}{L}\right),
\end{equation}
append step $o_{i^*,k_{i^*}}$ to the merged plan, and increment $k_{i^*}$.
This interleaves the children while preserving the internal order of every
child plan.

The decoded plan is then validated \emph{exactly}, not through the smooth
surrogates. For every shared ordering $a\prec b$, let $\pos(t)$ be the index of
the first step of the decoded plan whose task index is $t$; the constraint
holds iff both positions exist and $\pos(a)<\pos(b)$. A shared atom holds iff
it belongs to
\begin{equation}
s_0\;\cup\;\bigcup_{i=1}^{n_c}\bigl(\mathrm{final}_i\cup\mathrm{achieved}_i\bigr),
\end{equation}
where $\mathrm{final}_i$ are the local atoms true at the end of child $i$ and
$\mathrm{achieved}_i$ are the ground effect atoms of its operators. An equality
constraint is evaluated on its ground terms. If any constraint is still
violated, the merge reports an error with the residual and weight of every
offending constraint; there is no replanning at this stage.

\begin{proposition}[Merge soundness]
If a merge node returns successfully, then the decoded plan is a linearisation
of the child plans that satisfies every shared constraint of that node, and
its restriction to each child is exactly that child's plan.
\end{proposition}
\begin{proof}
The decoder appends the next unscheduled step of one child at a time, so the
projection onto each child is its original plan and its internal order is
preserved. Feasibility is established by the exact validator above, which
checks every shared ordering, atom and equality independently. A returned
merge has passed that check for all constraints.
\end{proof}

\begin{remark}
The continuous optimization is only a heuristic for \emph{choosing} the
interleaving. The smooth violation \eqref{eq:hinge} is positive at the
boundary, so feasibility is always decided by the exact test, never by the
surrogate value.
\end{remark}

\section{Hierarchical refinement}
\label{sec:refinement}

The abstract plan produced by the evaluation tree is a sequence of grounded
top-level tasks, each annotated with the method chosen by the abstract
operator. The refinement stage recursively replaces every compound task by a
method instance and every primitive task by an action, producing a complete
HTN plan.

\subsection{The refinement relation}
For a task $t$ and state $s$, let $\cR(t,s)$ be the finite set of pairs
$(\pi,s')$ such that $\pi$ is a primitive sequence executable from $s$ and
$s'$ is the resulting state. It is defined recursively:
\begin{itemize}[leftmargin=1.6em,itemsep=2pt]
\item If $t$ is a primitive action $a$ and $s\models\pre_a$, then
      $\cR(t,s)=\{(\langle a\rangle,\tau_a(s))\}$; otherwise
      $\cR(t,s)=\varnothing$.
\item If $t$ is compound, let $m\sigma$ range over the grounded method
      instances with head $t$ and $\phi_m\sigma$ true in $s$. For each such
      instance, refine the subtask network $TN_m\sigma$:
      \begin{itemize}[leftmargin=1.2em,itemsep=1pt]
      \item for a total order $t_1,\dots,t_n$, form the sequential product
            \[
            \cR(t_1\cdots t_n,s)
            =\Bigl\{(\pi_1\cdots\pi_n,s_n):
            (\pi_1,s_1)\in\cR(t_1,s),\dots,(\pi_n,s_n)\in\cR(t_n,s_{n-1})\Bigr\};
            \]
      \item for a partial order, enumerate the ready tasks (those whose
            predecessors are all scheduled), recurse on each ready task, and
            thread the state through the recursion. Equality and inequality
            constraints of the network are checked before expansion.
      \end{itemize}
      The results are wrapped into a method step
      $\langle t,m\sigma,\cdot\rangle$.
\end{itemize}

\subsection{Bounded search over the abstract plan}
Let the abstract plan be $s_1,\dots,s_m$ and let
$\cC_0=\{(\varepsilon,s_0)\}$. For $j=1,\dots,m$,
\begin{equation}
\cC_j=\operatorname{trunc}_{2048}
\Bigl(\bigcup_{(\rho,s)\in\cC_{j-1}}
\{\rho\cdot\pi:(\pi,s')\in\cR(s_j,s)\}\Bigr),
\end{equation}
where $\rho\cdot\pi$ concatenates the prefix $\rho$ with the new primitive
sequence $\pi$. If
$\cC_j=\varnothing$ the refinement fails. The bound is
$N_{\max}=2048$ partial refinements per abstract step, the recursion depth is
at most $32$, and method grounding is bounded by $B_g$. A path-local visited
set of $(\text{task},\text{state})$ pairs prunes branches that rediscover the
same subproblem; this is what stops recursive methods such as
\texttt{get\_to} through an intermediate location from cycling. The method
recorded in the abstract plan is tried first, then the remaining methods in
increasing order of subtask count.

After the last abstract step, the goal is evaluated on each candidate state.
The first candidate whose goal is satisfied (or absent) is returned; if none
satisfies the goal, the first candidate is returned as a fallback and the
caller can inspect \texttt{Refinement::goal}. The result is independently
checked by \texttt{validate\_refinement}, which re-executes method
preconditions, subtask orderings, primitive effects, task coverage and the
goal.

\subsection{Primitive method summaries}
A method
$m=\langle c(\vec{x}),\phi_m,\langle a_1\cdots a_n,\varnothing\rangle\rangle$
is summarised when every subtask $a_i$ is primitive, the network is totally
ordered, the method has no equality constraints, every action has exactly one
DNF precondition clause $\kappa_i$ and one unconditional effect alternative
$(A_i,D_i)$, and $A_i\cap D_i=\varnothing$.

The macro precondition is computed by backward regression. Set
$R=\varnothing$, $A=\varnothing$, $D=\varnothing$ and process
$i=n,n-1,\dots,1$:
\begin{enumerate}[leftmargin=1.6em,itemsep=1pt]
\item for every literal $\ell=(p,b)$ in $R$: if $p\in A_i$, require $b=1$
      (otherwise fail); if $p\in D_i$, require $b=0$ (otherwise fail);
      otherwise keep $\ell$ in $R$;
\item add the literals of $\kappa_i$ to $R$, failing on a polarity conflict;
\item for every $p\in A_i$ with $p\notin D$, insert $p$ into $A$;
\item for every $p\in D_i$ with $p\notin A$, insert $p$ into $D$.
\end{enumerate}
If no conflict occurred, the summary operator is
\begin{equation}
\pre(o)=\phi_m\;\wedge\;\bigwedge_{\ell\in R}\ell,
\qquad
\add(o)=A,
\qquad
\del(o)=D,
\qquad
\add(o)\mathrel{+}=\{\exec_t\}.
\label{eq:summary}
\end{equation}
This is exact under the stated restrictions: the subtasks are totally ordered,
deterministic and unconditional, so their composed effect is a single
add/delete pair and their required precondition is the regressed literal set.
Any method outside the fragment is left opaque.

\section{Optional learned initialization}

When the crate is compiled with the \texttt{neural} feature, the merger can be
warm-started by a heterogeneous graph-attention network instead of the
sequential start. The graph has one node per child plan, one node per shared
constraint and one global node, with $32$ features per node and a dense
incidence adjacency including self-loops. Three normalized GAT layers compute
\begin{equation}
\alpha_{ij}=\softmax_j\Bigl(\mathrm{LeakyReLU}\bigl(a^\top[Wh_i\,\|\,Wh_j]\bigr)\Bigr),
\qquad
h_i'=\sigma\Bigl(\sum_{j\in N(i)}\alpha_{ij}Wh_j\Bigr),
\end{equation}
and two heads predict the normalized start positions $q_i$ and the positive
penalty weights $w_c$. Training targets are produced by running the
deterministic merger from the sequential start and recording the optimized
parameters.

The predictions are only an initialization. They are accepted when their
shape matches the problem, all positions are finite and all weights are
strictly positive; otherwise the merger falls back to the sequential start.
The exact validator of Section~\ref{sec:merge} runs in all cases. The default
build has no Burn dependency and uses the deterministic sequential start.

A second optional model (\texttt{PSPR\_DECOMPOSITION\_MODEL}) can rank method
candidates and order ready subtasks during refinement. Its outputs are
accepted only when they form a permutation of the candidate set; otherwise the
baseline order is used. Both learned components are therefore guidance for
search ordering, not sources of logical soundness.

\section{Correctness, cost, and limitations}
\label{sec:analysis}

\subsection{What is guaranteed}
\begin{proposition}[Leaf soundness]
Every plan returned by a leaf search is executable in the leaf task and ends
in a state where all markers $e_t$ of the cluster are true. Each operator
corresponds to a grounded action or method instance of the original domain.
\end{proposition}

\begin{proposition}[Conditional completeness of the tree]
If every leaf is solved, every merge node succeeds, and the refinement stage
returns a candidate whose goal is satisfied, then the returned primitive plan
is a complete HTN plan of the initial task network and satisfies the state
goal.
\end{proposition}
\begin{proof}[Proof sketch]
The merge proposition gives, by induction over the schedule, a plan of every
subtree that satisfies all shared constraints attached to its nodes. The root
therefore satisfies all constraints of $TN_0$. Refinement replaces each
top-level task by a grounded method whose precondition holds in the state
threaded through the recursion, and its subtasks are recursively refined into
executable primitive actions. The goal check is performed on the final state,
and \texttt{validate\_refinement} repeats the check independently.
\end{proof}

\subsection{Sources of incompleteness}
The algorithm is deliberately satisficing and incomplete:
\begin{enumerate}[leftmargin=1.6em,itemsep=2pt]
\item \textbf{Abstraction gap.} $\alpha(P)$ is not a relaxation of $P$
      (Section~2), so an abstract plan can be unrefinable and a valid HTN plan
      can have no abstract counterpart.
\item \textbf{Dropped shared atoms.} Leaves cannot test or establish shared
      atoms. The merge can only reorder existing child plans; it cannot
      replan a child that needs a shared atom as a precondition.
\item \textbf{Bounded enumeration.} Grounding, DNF expansion, leaf search,
      refinement depth and refinement breadth are all capped.
\item \textbf{No recursive decomposition in the abstract tree.} Compound
      tasks in a leaf are opaque method operators; decomposition happens only
      in the refinement stage.
\item \textbf{Deterministic fragment.} Probabilistic formulas are rejected,
      and disjunctive effects are compiled as an angelic choice over
      deterministic variants.
\end{enumerate}

\subsection{Cost}
Let $n_1$ be the number of grounded tasks, $n_2$ the number of constraints,
$k$ the number of clusters, and $n_c$ the largest number of children of a
merge node.
\begin{itemize}[leftmargin=1.6em,itemsep=1pt]
\item Grounding and incidence construction are linear in the number of
      grounded tasks and their schema atoms, up to the budget $B_g$ and the
      DNF cap of $4096$.
\item Spectral clustering costs $O(n_1^3)$ per group for the eigendecomposition
      and $O(n_1\cdot\mathrm{nnz}(A))$ for the Gram matrix. The number of
      $k$-median iterations is bounded by the implementation.
\item The Gilpin ILP has $O(k^3)$ binary indicators and is NP-hard in general;
      it is solved by SCIP on the small cluster counts that occur in practice.
      The scheduling ILP has $O(k^2)$ binaries.
\item A leaf search is bounded by $10^5$ expansions and $5\times10^6$
      generated states.
\item A merge node costs $O(R\cdot S\cdot(n_c+|\cS_v|))$ for $R$ rounds of
      $S$ Adam steps, plus $O(L)$ for decoding. There are $k-1$ merge nodes,
      and every step of a child plan is decoded exactly once per node on the
      path to the root.
\end{itemize}

\subsection{Implementation status}
The pipeline described here is the one implemented by the current working
tree. The library test suite passes ($73$ tests at the time of writing); the
IPC integration tests are gated on an external HDDL benchmark checkout and the
\texttt{hddl\_analyzer} executable. The default build contains no neural code
and no Burn dependency. The evaluator reports a plan as \texttt{solved} when
every leaf is solved and the root merge succeeds; the state goal is reported
separately and is only enforced by the hierarchical path. The default planner
is therefore a fast solver of the task-network projection, while
\texttt{--hierarchical} (stages C+G) is the complete-HTN mode.

\begin{table}[h]
\centering
\small
\begin{tabular}{@{}ll@{}}
\toprule
Stage & Module \\
\midrule
HDDL parsing and JSON decoding & \texttt{src/parsing/} \\
Domain and problem translation & \texttt{src/solver/domain.rs}, \texttt{problem.rs} \\
Grounding, incidence graph, formula compiler & \texttt{evaluation\_tree/graph.rs}, \texttt{formula.rs} \\
Spectral clustering (BiSBM) & \texttt{partition/bisbm/} \\
Merge-level and scheduling ILPs & \texttt{partition/agglomeration/} \\
Tree construction and leaf evaluation & \texttt{evaluation\_tree/clustering.rs}, \texttt{leaf.rs} \\
Leaf search (FF, GBFS) & \texttt{fast\_downward/} \\
Relaxed SGPlan merger & \texttt{evaluation\_tree/merge/} \\
Hierarchical refinement (C+G) & \texttt{evaluation\_tree/hierarchical.rs} \\
Liftability gap detection & \texttt{evaluation\_tree/liftability.rs} \\
Independent validation & \texttt{evaluation\_tree/verify.rs} \\
Optional learned warm start & \texttt{evaluation\_tree/merge/neural.rs} \\
\bottomrule
\end{tabular}
\caption{Mapping from the stages of the algorithm to the implementation.}
\end{table}

\begin{thebibliography}{9}

\bibitem{sgplan}
Y.~Hsu, B.~Wah, R.~Huang, and Y.~Chen.
\newblock Constraint partitioning in penalty formulations for solving temporal
planning problems.
\newblock \emph{Artificial Intelligence}, 170(3):187--211, 2006.

\bibitem{gilpin}
S.~Gilpin, T.~Sandholm, and J.~S{\o}rensen.
\newblock Formalizing hierarchical clustering as integer linear programming.
\newblock In \emph{Proceedings of AAAI}, 2013.

\bibitem{braun}
F.~Braun.
\newblock Strong consistency guarantees for clustering high-dimensional
bipartite graphs with the spectral method.
\newblock arXiv:2304.06903, 2023.

\bibitem{helmert}
M.~Helmert.
\newblock The Fast Downward planning system.
\newblock \emph{Journal of Artificial Intelligence Research}, 26:191--246, 2006.

\bibitem{hoffmann}
J.~Hoffmann and B.~Nebel.
\newblock The FF planning system: fast plan generation through heuristic
search.
\newblock \emph{Journal of Artificial Intelligence Research}, 14:253--302, 2001.

\bibitem{kingma}
D.~P. Kingma and J.~Ba.
\newblock Adam: a method for stochastic optimization.
\newblock In \emph{Proceedings of ICLR}, 2015.

\end{thebibliography}

\end{document}
